\subsection{Inside Job (2010)}

El documental examina las causas y dinámicas institucionales que condujeron a la crisis financiera global de 2008, enfocándose especialmente en el papel de los grandes bancos de inversión, las agencias calificadoras, los reguladores y la academia económica. A diferencia de narrativas más generales sobre el colapso financiero, la película presenta testimonios directos de economistas, funcionarios y actores del sector financiero para mostrar cómo ciertas decisiones y estructuras regulatorias contribuyeron a amplificar los riesgos dentro del sistema. En particular, se describe el crecimiento de instrumentos financieros complejos, como los CDOs y los credit default swaps, así como la forma en que estos productos se integraron en los mercados globales sin un marco regulatorio suficientemente robusto.

Uno de los aspectos más llamativos del documental es la evidencia presentada sobre cómo algunas instituciones financieras diseñaron y comercializaron productos estructurados basados en hipotecas subprime que recibían calificaciones crediticias altas, como “AAA”, a pesar de que contenían activos de calidad dudosa. Estos instrumentos fueron vendidos a inversionistas institucionales, incluyendo fondos de pensiones, que confiaban en las evaluaciones de riesgo proporcionadas por las agencias calificadoras. Resulta particularmente sorprendente que instituciones como Goldman Sachs continuaran operando con normalidad después de la crisis, a pesar de las prácticas documentadas en torno a la estructuración y comercialización de estos productos. Este hecho plantea interrogantes sobre la relación entre responsabilidad institucional, regulación financiera y la capacidad del sistema para sancionar conductas que pueden generar consecuencias sistémicas significativas.

El documental también señala la importancia de la desregulación financiera previa a la crisis, especialmente en el mercado de derivados. Durante años, muchos contratos derivados se negociaron en mercados extrabursátiles con niveles relativamente bajos de supervisión. En este contexto surgieron instrumentos como los credit default swaps, que permitían asegurar o apostar contra el riesgo de incumplimiento de determinados activos financieros. En algunos casos, estos contratos podían adquirirse incluso sin poseer el activo subyacente, lo que permitía a los participantes del mercado especular sobre el deterioro de ciertos instrumentos financieros. Esta estructura plantea una pregunta fundamental: ¿cómo fue posible asegurar riesgos que, en conjunto, superaban ampliamente la capacidad real de pago de quienes vendían esos contratos? El documental sugiere que la combinación de innovación financiera, incentivos económicos y debilidad regulatoria permitió que este sistema creciera hasta alcanzar niveles insostenibles.

No necesariamente se trata de afirmar que las empresas o los mercados actúen con una intención intrínsecamente negativa. En muchos casos, los agentes económicos simplemente responden a los incentivos que enfrentan. Las corporaciones financieras, como otras organizaciones económicas, tienden a maximizar beneficios dentro de las reglas existentes. Sin embargo, cuando estas instituciones ocupan una posición central dentro del sistema económico, sus decisiones pueden afectar a toda la sociedad. Por esta razón, la existencia de regulaciones adecuadas resulta fundamental para establecer límites claros dentro de los cuales las instituciones puedan operar.

Desde esta perspectiva, la regulación financiera no debe entenderse como una restricción arbitraria a la actividad económica, sino como un mecanismo necesario para alinear los incentivos privados con la estabilidad del sistema en su conjunto. Si las sanciones o multas asociadas a ciertas prácticas son menores que los beneficios obtenidos al realizarlas, las instituciones pueden encontrar racional continuar con ese tipo de conductas. En consecuencia, el diseño regulatorio debe asegurar que el costo esperado de infringir las normas supere las ganancias potenciales derivadas de hacerlo. Este principio es particularmente relevante en el caso de las instituciones financieras sistémicamente importantes, cuyos riesgos pueden propagarse rápidamente al resto de la economía.

Más que atribuir el colapso a individuos aislados, el documental sugiere que el problema reside en una estructura de incentivos que permitió que prácticas riesgosas se desarrollaran durante años sin una supervisión adecuada. En este contexto, el fortalecimiento de la regulación financiera aparece como una condición necesaria para mantener la estabilidad del sistema y evitar que decisiones racionales desde el punto de vista individual produzcan consecuencias negativas a nivel colectivo.
